\documentclass[12pt]{article}
\usepackage[utf8]{inputenc}
\usepackage[T1]{fontenc}
\usepackage{hyperref}
\title{Vaghri Language Resources for K-12 Curriculum Development}
\author{Compiled by Claude Sonnet 5 (claude-sonnet-5) for Andr\'as Kornai}
\date{}
\begin{document}
\maketitle

\section{Introduction}
Vaghri (vgr) is an Indo-Aryan language of Sindh, Pakistan, in Umerkot and Tharparkar districts. It is closely related to Gujarati. Its identity relative to Wagdi and Bauria is unclear, and speaker counts range from 4,800 to 9,000. Little exists, and what was found is mostly missionary material.

\section{Primary Educational Resources}
None found.

\section{Linguistic and Reference Materials}
Overview pages only: \url{https://en.wikipedia.org/wiki/Vaghri_language} and \url{https://en.wikipedia.org/wiki/Vaghri_people}

\section{Religious and Community Resources}
Scripture Earth lists a Vaghri New Testament: \url{https://scriptureearth.org/00eng.php?iso=vgr}

\section{Digital and Audio Resources}
Global Recordings Network has audio recordings in Vaghri: \url{https://globalrecordings.net/en/language/vgr}

\section{Curriculum Development Context}
Vaghri is not known to be taught in schools. No curriculum, primer or dictionary was found, so school materials would need to be created from scratch with the community.
\end{document}